% !TEX root = ../../hw07.tex
% Source: author's supplied text in chat, 2026-11-06, Q10.
% Preserve logarithmic transformation and L'Hopital calculation.
% Every approach to zero is from the positive domain of the real power.
\begin{proof}[Solution]
  We take the limit from the right because the
  real-power expression here is defined for $x>0$.
  Applying the same transformation from
  \autoref{ex:hw07:self-power-derivative} gives
  \[
    \lim_{x\to0^+}x^x
      =\lim_{x\to0^+}e^{\ln(x^x)}
      =\lim_{x\to0^+}e^{x\ln x}.
  \]
  Since the exponential function is continuous everywhere on $\R$,
  we may interchange it with the limit, provided the inner limit
  exists and is finite,
  \[
    \lim_{x\to0^+}e^{x\ln x}
      =e^{\lim_{x\to0^+}x\ln x}.
  \]
  We now evaluate the exponent. Rewriting the product as a quotient
  gives the indeterminate form $-\infty/\infty$, allowing us to
  apply L'H\^opital's rule.
  More precisely, apply the $+\infty/+\infty$
  version to $(-\ln x)/(1/x)$ and then restore the minus sign.
  The functions are differentiable for $x>0$, the denominator
  derivative $-1/x^2$ is nonzero, and the derivative ratio below
  has a limit.
  % Keep the exponent calculation with its concluding substitution.
  \postdisplaypenalty=10000
  \[
    \begin{aligned}
      \lim_{x\to0^+}x\ln x
        &=\lim_{x\to0^+}\frac{\ln x}{1/x}\\
        &\overset{\mathrm{L'H}}{=}
          \lim_{x\to0^+}\frac{1/x}{-1/x^2}\\
        &=\lim_{x\to0^+}(-x)\\
        &=0.
    \end{aligned}
  \]
  Therefore, substituting this result into the original expression,
  \[
    \lim_{x\to0^+}x^x
      =e^{\lim_{x\to0^+}x\ln x}=e^0=1.
  \]
\end{proof}
